% Converted from README.md
\section{Pisanje testabilnog koda}\label{pisanje-testabilnog-koda}

Da bismo efikasno pisali testove jedinica koda, neophodno je da kod pišemo tako da bude pogodan za testiranje, ali i da softver razvijamo tako da je moguće paralelno pisati testove. Jedan način da se obezbedi takav način rada je razvoj vođen testovima (engl. \emph{Test-Driven-Development}, skr. \emph{TDD}). U situacijama kada već imamo dostupan izvorni kod i treba da napišemo testove jedinica koda, treba znati kako pristupiti pisanju testova i kako refaktorisati kod tako da bude pogodan za testiranje. Nekada funkcije koje treba da testiramo nemaju deterministički izlaz tako da ih je nemoguće testirati u izvornom obliku.

U ovom poglavlju će biti reči o tehnikama za pisanje testabilnog i skalabilnog koda. Primeri na kojima će ove tehnike biti prikazane su pisani u raznim programskim jezicima ali su koncepti koje prikazuju univerzalni, kao npr. inverzija zavisnosti, zamena implementacije itd.

\subimport{01_cpp_calculator/}{index.tex}

\subimport{02_java_gol/}{index.tex}

\subimport{03_csharp_service/}{index.tex}
